% =========================================================
% Generalized Heine-Borel Theorem
% Source: volume-v/topology/point-set-topology/notes/
%
% Dependencies:
%   notes-metric-spaces (completeness, totally bounded)
%   notes-topology (compactness via open covers)
%   notes-sequential-characterizations (sequential compactness)
% =========================================================

% ---------------------------------------------------------
\section{Generalized Heine-Borel Theorem}
\label{sec:generalized-heine-borel}
% ---------------------------------------------------------

\begin{tcolorbox}[colback=gray!6, colframe=gray!40, arc=2pt,
  left=6pt, right=6pt, top=4pt, bottom=4pt,
  title={\small\textbf{Generalized Heine-Borel --- Quick Reference}},
  fonttitle=\small\bfseries]
\small
\begin{tabular}{l l l l}
\toprule
\textbf{Label} & \textbf{Statement} & \textbf{Proof method} & \textbf{Detail} \\
\midrule
T\ref{thm:gen-heine-borel}
  & Compact $\iff$ complete $+$ totally bounded
  & Sequential compactness
  & \hyperref[thm:gen-heine-borel]{$\downarrow$} \\[4pt]
C\ref{cor:classical-heine-borel}
  & In $\mathbb{R}^n$: compact $\iff$ closed $+$ bounded
  & Totally bounded $\iff$ bounded in $\mathbb{R}^n$
  & \hyperref[cor:classical-heine-borel]{$\downarrow$} \\[4pt]
\bottomrule
\end{tabular}
\end{tcolorbox}

\paragraph{Motivation.}
The classical Heine-Borel theorem says compact subsets of $\mathbb{R}^n$ are
exactly the closed and bounded ones.
Once you know what compactness and total boundedness mean in general metric
spaces, a natural question arises: which condition is really doing the work?
The answer is total boundedness, not mere boundedness.
The Generalized Heine-Borel theorem isolates the two ingredients ---
completeness and total boundedness --- that together force every sequence
to have a convergent subsequence.

\begin{tcolorbox}[colback=thmbox, colframe=thmborder, arc=2pt,
  left=6pt, right=6pt, top=4pt, bottom=4pt,
  title={\small\textbf{Theorem (Generalized Heine-Borel --- Pugh Theorem 65)}},
  fonttitle=\small\bfseries]
\begin{theorem}[Generalized Heine-Borel]
\label{thm:gen-heine-borel}
Let $(X,d)$ be a complete metric space and $A \subseteq X$.
Then $A$ is compact if and only if $A$ is closed and totally bounded.
\end{theorem}
\end{tcolorbox}

\begin{remark*}[Fully quantified form]
\[
\forall (X,d)\ \text{complete metric space},\;
\forall A \subseteq X:\quad
A\ \text{compact}
\iff
\bigl(A\ \text{closed}\bigr)
\land
\bigl(\forall\,\varepsilon>0,\;
\exists\, n\in\mathbb{N},\;
\exists\, x_1,\ldots,x_n \in X,\;
A \subseteq \bigcup_{k=1}^n B_\varepsilon(x_k)\bigr).
\]
\end{remark*}

\begin{remark*}[Proof strategy]
($\Rightarrow$) Every compact set is closed (as a closed subspace of a complete
space it is complete in its own right) and totally bounded (the open cover by
$\varepsilon$-balls has a finite subcover).

($\Leftarrow$) Let $(a_n)$ be any sequence in $A$.
Since $A$ is totally bounded, Proposition~\ref{prop:totally-bounded-Cauchy}
(in the metric-spaces chapter) produces a Cauchy subsequence $(b_n)$.
Since $X$ is complete, $(b_n)$ converges to some $p \in X$.
Since $A$ is closed, $p \in A$.
So every sequence in $A$ has a subsequence converging in $A$ ---
which is sequential compactness, hence compactness.
The equivalence of sequential compactness and open-cover compactness
in metric spaces is established in the compactness chapter.
\end{remark*}

\begin{tcolorbox}[colback=corbox, colframe=corborder, arc=2pt,
  left=6pt, right=6pt, top=4pt, bottom=4pt,
  title={\small\textbf{Corollary (Classical Heine-Borel in $\mathbb{R}^n$)}},
  fonttitle=\small\bfseries]
\begin{corollary}[Classical Heine-Borel]
\label{cor:classical-heine-borel}
A subset $K \subseteq \mathbb{R}^n$ is compact if and only if
$K$ is closed and bounded.
\end{corollary}
\end{tcolorbox}

\begin{remark*}[Why classical and generalized coincide in $\mathbb{R}^n$]
In $\mathbb{R}^n$, bounded and totally bounded are equivalent
(Proposition~\ref{prop:totally-bounded-Rn} in the metric-spaces chapter).
So in $\mathbb{R}^n$, the Generalized Heine-Borel theorem reads exactly as
the classical one: compact $\iff$ closed $+$ bounded.
The generalized version reveals that it is total boundedness, not bare
boundedness, that is the correct condition in general metric spaces.
\end{remark*}

\begin{remark*}[Why Heine-Borel fails in infinite dimensions]
In $\ell^2$ (the space of square-summable sequences), the closed unit ball
$B = \{x : \|x\|_2 \le 1\}$ is closed and bounded but not compact.
The standard basis vectors $e_n$ satisfy $\|e_m - e_n\|_2 = \sqrt{2}$ for
$m \ne n$, so no subsequence is Cauchy, and $B$ is not sequentially compact.
The unit ball in $\ell^2$ is bounded but not totally bounded, which is the
correct diagnosis: total boundedness, not boundedness, is what fails.

This failure is fundamental to infinite-dimensional analysis and motivates
the study of compact operators in functional analysis.
\end{remark*}

\begin{remark*}[The role of completeness]
The hypothesis that $X$ is complete is essential.
A totally bounded but incomplete space can fail to be compact:
the open interval $(0,1)$ with the usual metric is totally bounded
(it sits inside $[0,1]$ which is compact) but is not complete,
and it is not compact (the sequence $1/n$ has no limit in $(0,1)$).
\end{remark*}
